% §19.5 Composition, Spectra, and the Algebra
%
% Monograph register: gather the payoff. (a) Axiom 1 <=> product operator;
% parsimony = factorization (recap from §19.4). (b) Composition = operator
% product (eigenvalues multiply), giving prod(1-eta_i). (c) Spectral/Walsh view
% of parity: symmetric q-flip channel [[1-q,q],[q,1-q]] has eigenvalues 1 and
% 1-2q; 1-2q is simultaneously the +/-1 bias E[(-1)^flip] and the Walsh
% coefficient, so parity of n independent bits errs with prob
% (1/2)(1 - prod_i(1-2q_i)) = (1/2)(1 - product of nontrivial eigenvalues).
% Note: this is the syndrome-bit identity of the quantum-readout companion and
% the classical noisy-formula tradition (von Neumann/Pippenger); three names
% (eigenvalue / Walsh coefficient / +/-1 bias) for one quantity. (d) Set
% operations in the tensor algebra: union = OR on bits, intersection = AND; the
% intersection FPR is the three-region weighted rate of eq:intersection-fpr
% (eps_A eps_B is only the sparse corner), so the operator view does not erase
% that subtlety. Show operator statement ALONGSIDE the set-theoretic one.
% Cross-refs sec:operator-sets-tensored, prop:set-operator-is-axiom-1,
% def:set-noise-operator, prop:operator-map-composition, thm:composition,
% eq:composition-correctness, sec:other-set-ops, prop:intersection-fpr-fnr,
% eq:intersection-fpr, eq:union-fpr, eq:union-fnr, def:union-of-sets,
% def:intersection-of-sets. DEFINED: sec:operator-composition-spectra,
% def:walsh-eigenvalue, prop:parity-spectral.

\section{Composition, Spectra, and the Algebra}
\label{sec:operator-composition-spectra}

With maps as matrices and sets as tensored channels in place, three things the
set-theoretic chapters proved separately become facets of one algebraic
picture. This section collects them, in each case setting the operator statement
beside the predicate statement it matches.

\subsection*{Independence is a product, and parsimony is the factoring}

The first facet is a recap in one sentence.
\Cref{prop:set-operator-is-axiom-1} identified element-wise independence with a
product operator: the set view's ``the errors at distinct elements are
independent'' is the operator view's ``$Q_{\mathrm{set}} = \bigotimes_x Q_x$.''
The parsimony that followed, $2N$ parameters in place of $2^N (2^N - 1)$, is the
arithmetic of a factored operator, whose parameter budget is the sum of its
factors' budgets. Independence, factorization, and parsimony are three readings
of the single equation $Q_{\mathrm{set}} = \bigotimes_x Q_x$.

\subsection*{Composition multiplies eigenvalues}

The second facet sharpens \Cref{prop:operator-map-composition}. Serial
composition of Bernoulli maps is the matrix product $Q_g Q_f$, and a matrix
product does more than compose: when the factors are simultaneously
diagonalizable, it multiplies eigenvalues. Each one-sided survival probability
$1 - \eta_i$ is an eigenvalue of stage $i$ on the eigenvector that carries the
truth through unflipped, and the chain's survival probability
$\prod_i (1 - \eta_i)$ is the product of those eigenvalues. The composition
theorem's
\[
  \eta_{\mathrm{total}} \;=\; 1 - \prod_{i=1}^{k}(1 - \eta_i)
  \tag{\Cref{eq:composition-correctness}}
\]
is therefore a spectral statement: one minus a product of eigenvalues. The set
view obtained it by induction over chained gates; the operator view reads it off
the spectrum of a product. The next facet makes the spectral reading do work the
combinatorial form does not display.

\subsection*{Parity: one quantity under three names}

The cleanest place the spectrum earns its keep is the parity of noisy bits. Take
the symmetric single-bit flip channel
\[
  Q(q) \;=\;
  \begin{pmatrix}
    1 - q & q \\
    q & 1 - q
  \end{pmatrix},
\]
which flips a bit with probability $q$.

\begin{definition}[The nontrivial eigenvalue of a flip channel]
\label{def:walsh-eigenvalue}
The symmetric flip channel $Q(q)$ has eigenvectors $(1, 1)^{\mathsf{T}}$ and
$(1, -1)^{\mathsf{T}}$ with eigenvalues $1$ and $1 - 2q$. The first eigenvalue,
on the all-ones vector, is the stochasticity of the channel and carries no
information; the second, $\lambda(q) = 1 - 2q$, is the \emph{nontrivial
eigenvalue}, the factor by which the channel contracts the sign coordinate
$(1, -1)^{\mathsf{T}}$.
\end{definition}

The number $\lambda(q) = 1 - 2q$ has two further identities that coincide with
the eigenvalue. Read the bit through the sign map $b \mapsto (-1)^{b}$, sending
$\{0, 1\}$ to $\{+1, -1\}$. The expected sign after one flip is
\[
  \Expect{(-1)^{\text{flip}}}
  = (1 - q)(+1) + q(-1)
  = 1 - 2q,
\]
the \emph{$\pm 1$ bias} of the channel. And $1 - 2q$ is the order-one
\emph{Walsh (Fourier) coefficient} of the channel on the Boolean cube, the
coefficient of the single-bit character $\chi(b) = (-1)^{b}$ under the
expansion. Eigenvalue, $\pm 1$ bias, and Walsh coefficient are three names for
$1 - 2q$, and the names matter because each makes the parity formula obvious
from a different direction.

\begin{proposition}[Spectral error rate of a noisy parity]
\label{prop:parity-spectral}
Let $b_1, \ldots, b_n$ be bits, each flipped independently by a symmetric
channel of rate $q_i$, and let the output be the parity (XOR) $b_1 \oplus \cdots
\oplus b_n$. The parity is wrong with probability
\[
  \Prob{\text{parity flipped}}
  \;=\; \frac{1}{2}\Big(1 - \prod_{i=1}^{n}(1 - 2 q_i)\Big)
  \;=\; \frac{1}{2}\Big(1 - \prod_{i=1}^{n} \lambda(q_i)\Big),
\]
half of one minus the product of the channels' nontrivial eigenvalues.
\end{proposition}

\begin{proof}
The parity flips iff an odd number of the $n$ bits flip. Pass to signs: the
output sign is $\prod_i (-1)^{b_i'}$ where $b_i'$ is the flipped bit, and the
parity is wrong exactly when the product of flip signs is $-1$. By independence
the expected product of flip signs factors,
\[
  \Expect{\textstyle\prod_i (-1)^{\text{flip}_i}}
  = \prod_{i=1}^{n} \Expect{(-1)^{\text{flip}_i}}
  = \prod_{i=1}^{n}(1 - 2 q_i),
\]
each factor the $\pm 1$ bias of one channel. Writing $P$ for the probability the
parity is flipped, the expected sign of the error is $(+1)(1 - P) + (-1)P = 1 -
2P$, so $1 - 2P = \prod_i (1 - 2q_i)$ and $P = \tfrac{1}{2}(1 - \prod_i (1 -
2q_i))$. Each $1 - 2q_i$ is the nontrivial eigenvalue $\lambda(q_i)$ of
\Cref{def:walsh-eigenvalue}, giving the spectral form.
\end{proof}

The product of $\pm 1$ biases is the multiplicativity of the order-one Walsh
coefficient under XOR, and it is the same multiplicativity as eigenvalues under
a tensor product, since the parity error operator is built from the per-bit
channels by the Kronecker product of \Cref{def:set-noise-operator} restricted to
the sign sector. This identity is the syndrome-bit calculation used in the
quantum-readout setting, where a measured parity (a stabilizer syndrome) inherits
$\prod_i (1 - 2q_i)$ from its constituent single-qubit readout errors, and it is
the classical noisy-formula tradition begun for redundant computation by von
Neumann and analyzed for bounded-depth circuits by Pippenger, where the
$\pm 1$-bias product is the standard accounting for error through XOR gates. The
operator view unifies them: eigenvalue, Walsh coefficient, and $\pm 1$ bias are
one quantity, and a parity of independent noisy bits contracts the sign sector by
the product of that quantity over the bits.

\subsection*{Set operations in the tensor algebra}

The last facet returns to the set algebra of \Cref{ch:set-composition} and reads
it on configurations. Union and intersection of approximate sets are bitwise
logic on membership configurations: the union $\ASet{A} \cup \ASet{B}$
(\Cref{def:union-of-sets}) is the OR of the two membership bits at each element,
and the intersection $\ASet{A} \cap \ASet{B}$ (\Cref{def:intersection-of-sets})
is the AND. On the tensor space $\mathcal{V}$ both are elementwise (per-factor)
Boolean operations on the reported configurations, the predicate-level
operations of chapter~6 written one factor at a time.

The operator view does not, however, erase the subtlety chapter~6 was careful
about. The intersection's false-positive rate is not the bare product $\fprate_A
\fprate_B$; it is the partition-weighted three-region sum
\[
  \fprate_\cap
  = w_1 (1 - \fnrate_A)\fprate_B
  + w_2 \fprate_A (1 - \fnrate_B)
  + w_3 \fprate_A \fprate_B
  \tag{\Cref{eq:intersection-fpr}}
\]
of \Cref{prop:intersection-fpr-fnr}, where $w_1, w_2, w_3$ weight the regions
$A \setminus B$, $B \setminus A$, and outside both. The product $\fprate_A
\fprate_B$ is only the third term, the both-outside corner, exact in the sparse
limit where almost every non-member of $A \cap B$ lies outside both operands.
The reason survives translation: a false positive of the intersection requires
both operands to report membership at an element where the intersection is
absent, and that element's true status differs across the three regions, so the
per-element AND consults different entries of $Q_A$ and $Q_B$ depending on the
region. The elementwise tensor algebra computes the rate correctly only when it
is fed the right per-element truth, which is the same bookkeeping the predicate
derivation does by hand. The union is the De~Morgan dual, its false-negative
rate the three-region sum \Cref{eq:union-fnr} and its false-positive rate the
clean product $1 - (1 - \fprate_A)(1 - \fprate_B)$ of \Cref{eq:union-fpr}. The
operator picture makes the logic per-factor and the composition multiplicative;
it leaves the region accounting exactly where chapter~6 put it.
